\documentclass[11pt]{article}
\usepackage[a4paper,margin=25mm]{geometry}
\usepackage[T1]{fontenc}
\usepackage{lmodern,amsmath,amssymb,amsthm,booktabs,array,longtable}
\usepackage[colorlinks=true,urlcolor=blue,citecolor=blue,linkcolor=blue]{hyperref}
\newtheorem{theorem}{Theorem}
\newtheorem{corollary}[theorem]{Corollary}
\newcommand{\coeff}[2]{[#1^{#2}]}
\newcommand{\N}{\mathbb{N}}
\title{Kurkov's convolution-row conjectures over arbitrary semirings}
\author{Haobo Ma\\\small ChronoAI Pte Ltd, Singapore
\and Wenlin Zhang\\\small National University of Singapore\\\small ChronoAI Pte Ltd, Singapore}
\date{13 November 2026}
\begin{document}
\maketitle
\begin{abstract}
Mikhail Kurkov's FORMULA conjectures of 5 April 2025 in OEIS A340991,
revision 24, and A357583, revision 10 (both dated 6 April 2025), propose
the same shifted-array recursion for the prime and Bell convolution
triangles and for arbitrary coefficient sequences. We prove the recursion
over every unital semiring, including noncommutative ones, with the stated
order of multiplication. A shifted-tail coefficient formula determines
every array column. It also gives the Bell specialization, finite-support
recurrences, and the factorial generating function conjectured by Bala;
the latter and its continued fraction are deductions of classical INVERT
and Euler identities.
\end{abstract}

\section{The convolution triangle and the conjecture}
Kurkov conjectures that ``row polynomials are $xR(n,1)$'' for the prime
and Bell convolution triangles \cite{a340991,bell}. His accompanying
observation is that ``The same recursion seems to work for self-convolution
of any other sequence.'' We prove both conjectures by expressing each
array column using the corresponding tail of the coefficient sequence.

Let $K$ be a unital semiring: addition is a commutative monoid, multiplication
is an associative monoid, multiplication distributes over addition, and
$0$ annihilates products. Multiplication need not be commutative. Let
$a_1,a_2,\ldots$ be elements of $K$. The variables $z$ and $x$ are central,
so they commute with each other and with every coefficient. For a formal
power series $F$, write $[z^n]F$ for its coefficient of degree $n$. Set
\begin{equation}\label{eq:rows}
 A(z)=\sum_{j\geq1}a_jz^j,\qquad
 T(n,r)=[z^n]A(z)^r,\qquad
 P_n(x)=\sum_{r=0}^{n}T(n,r)x^r .
\end{equation}
Here $n,r\in\N=\{0,1,2,\ldots\}$, $A^0=1$, and $P_0(x)=1$.
Since $A$ has constant coefficient zero, $T(n,r)=0$ if $r>n$.
For commutative coefficients these $T(n,r)$ are the ordinary partial Bell,
or De Moivre, polynomials: see O'Sullivan \cite[Definition 1.1 and
equations (2.4), (2.6)]{osullivan}. In particular, the ordered-composition
expansion of $A^r$ is classical; our argument identifies it with Kurkov's
shifted array.

For $n,k\geq1$, define polynomials $R(n,k)$ recursively by
\begin{equation}\label{eq:recursion}
 R(1,k)=a_k,\qquad
 R(n,k)=R(n-1,k+1)+xR(n-1,1)a_k\quad(n>1).
\end{equation}
Taking $a_j$ to be the $j$th prime gives the exact recursion conjectured by
Mikhail Kurkov on 5 April 2025 in OEIS A340991 \cite{a340991}; taking $a_j$
to be the $j$th Bell number gives his identical conjecture in A357583
\cite{bell}. Both entries propose the same rule for any sequence. Keeping
the factor $a_k$ on the right in \eqref{eq:recursion} makes that assertion
valid even when coefficients do not commute.

\begin{theorem}\label{thm:main}
For every sequence over every unital semiring and every $n,k\geq1$,
\begin{equation}\label{eq:tailformula}
 R(n,k)=\sum_{r=0}^{n-1}x^r[z^{n-1}]A(z)^rB_k(z),
 \qquad B_k(z)=\sum_{j\geq0}a_{k+j}z^j.
\end{equation}
Consequently, for every $n\geq1$,
\begin{equation}\label{eq:identity}
 P_n(x)=xR(n,1).
\end{equation}
\end{theorem}

\begin{proof}
All series calculations are formal. Work in $K[x][[z]]$ and put
\[
 S(z,x)=\sum_{r\geq0}x^rA(z)^r,\qquad C_k(z,x)=S(z,x)B_k(z).
\]
These expressions are well defined: in degree $n$ in $z$, only the
finitely many indices $r\leq n$ contribute, so every coefficient is a
polynomial in $x$. The elementary identities
\begin{equation}\label{eq:splitting}
 B_k=a_k+zB_{k+1},\qquad A=zB_1,\qquad S=1+xSA
\end{equation}
follow coefficientwise. For the last identity, multiply each term of
$S$ on the right by $A$ and shift the exponent $r$ by one; this uses
associativity, not commutativity of coefficients. The order of the two
factors in $C_k=SB_k$ is important. Equations \eqref{eq:splitting} give
\begin{align*}
 C_k
 &=Sa_k+zSB_{k+1}\\
 &=a_k+xSAa_k+zC_{k+1}\\
 &=a_k+zC_{k+1}+xzC_1a_k.
\end{align*}
The constant coefficient of $C_k$ is $a_k$. Extracting the coefficient
of $z^{n-1}$ for $n>1$ shows that the polynomials
$[z^{n-1}]C_k$ obey exactly \eqref{eq:recursion}. Induction on $n$,
simultaneously for every positive $k$, therefore gives
$R(n,k)=[z^{n-1}]C_k$. Expanding $S$ proves
\eqref{eq:tailformula}; terms with $r>n-1$ vanish.

Also $S=1+xzC_1$ by \eqref{eq:splitting}. The coefficient of $z^n$ in
$S$ is $P_n(x)$ by \eqref{eq:rows}. For $n\geq1$ the constant series
$1$ contributes zero, so $P_n(x)=x[z^{n-1}]C_1=xR(n,1)$.
This proves \eqref{eq:identity}. The zeroth row is $P_0=1$ and is not
part of the positive-row assertion.
\end{proof}

For example, over a noncommutative semiring the fourth row is
\begin{align*}
 P_4(x)={}&a_4x+(a_1a_3+a_2^2+a_3a_1)x^2\\
 &+(a_1^2a_2+a_1a_2a_1+a_2a_1^2)x^3+a_1^4x^4 .
\end{align*}
For $a_1,a_2,a_3,a_4=2,3,5,7$ this becomes
$7x+29x^2+36x^3+16x^4$, the fourth positive row displayed in
A340991.

\section{Consequences and source coverage}
The same coefficient expansion gives useful specializations without the
auxiliary array. In the following display only, assume $K$ is a
commutative ring, so the indicated formal inverses exist.

\begin{corollary}\label{cor:sums}
For $n\geq0$,
\begin{align}
 P_n(1)&=[z^n]\frac{1}{1-A(z)},&
 P_n(-1)&=[z^n]\frac{1}{1+A(z)},\label{eq:sums}\\
 \sum_{r=0}^n rT(n,r)&=[z^n]\frac{A(z)}{(1-A(z))^2}.
 &&\label{eq:weighted}
\end{align}
For a sequence over any unital semiring, $T(n,n)=a_1^n$ and
\begin{equation}\label{eq:twosided}
 P_n(x)=x\sum_{j=1}^n P_{n-j}(x)a_j
       =x\sum_{j=1}^n a_jP_{n-j}(x)\quad(n\geq1).
\end{equation}
If $a_j=0$ for $j>d$, both sums stop at $\min(n,d)$.
\end{corollary}
\begin{proof}
Evaluate $S=\sum_{r\geq0}x^rA^r$ at $x=1$ and $x=-1$ for
\eqref{eq:sums}. Differentiate this identity formally with respect to
the central variable $x$, evaluate at $x=1$, and use
$\partial_x(1-xA)^{-1}=A(1-xA)^{-2}$ for \eqref{eq:weighted}.
All operations are coefficientwise finite. In $A^n$, the only product
of $n$ positive degrees summing to $n$ uses degree one each time,
giving $T(n,n)=a_1^n$, including $T(0,0)=1$. Finally, taking
$[z^n]$ in $S=1+xSA=1+xAS$ gives the two recurrences. The second
series identity holds because $A$ commutes with each of its own powers,
even if distinct coefficients do not commute. Zero coefficients beyond
$d$ remove the remaining summands.
\end{proof}

For primes, these are the established row-sum identities A030017$(n+1)$,
A030018$(n)$ and A030281$(n)$, respectively \cite{sequences}.
The first identity is the classical INVERT transform of Bernstein and
Sloane \cite[\S1, INVERT]{bernstein}; see also \cite{transforms}.

\begin{corollary}\label{cor:bell}
Let $b_j=\operatorname{Bell}(j)$ for $j\geq0$ and
$B(z)=\sum_{j\geq1}b_jz^j$. Kurkov's recursion with $R(1,k)=b_k$
gives every positive row of A357583. Its row sums $u_n$ satisfy
\[
 u_0=1,\qquad u_n=\sum_{j=1}^n b_j u_{n-j},\qquad
 \sum_{n\geq0}u_nz^n=\frac1{1-B(z)}.
\]
Its central terms are $[z^{2n}]B(z)^n$.
\end{corollary}
\begin{proof}
Substitute $a_j=b_j$ into Theorem~\ref{thm:main} and
Corollary~\ref{cor:sums}, and then take $x=1$ in
\eqref{eq:twosided}. Taking column $n$ of row $2n$ gives the last formula.
The row sums and central terms are the defining identities of A129247
and A357584, respectively \cite{bell}; they are not additional conjectures.
\end{proof}

\begin{corollary}\label{cor:factorial}
Set $a_j=j!$ for $j\geq1$ and $F(z)=\sum_{j\geq0}j!z^j$. The
row polynomials of A090238 have the bivariate generating function
\begin{equation}\label{eq:bala}
 G(z,x):=\sum_{n\geq0}P_n(x)z^n=\frac1{1+x-xF(z)}.
\end{equation}
Writing $H(z)=\cfrac{z}{1-\cfrac{2z}{1-\cfrac{2z}{1-\cfrac{3z}{1-\cfrac{3z}{\ddots}}}}}$,
the same series is
\begin{equation}\label{eq:balacf}
 G(z,x)=1-x+\frac{x}{1-\dfrac{z}{1-xz-H(z)}}.
\end{equation}
These are the generating function and continued fraction proposed by
Peter Bala in the FORMULA section of A090238 on 27 May 2017
\cite{factorial}; his variables $x,t$ are $z,x$ here.
\end{corollary}
\begin{proof}
Here $A=F-1$, so the geometric series $S=(1-xA)^{-1}$ proves
\eqref{eq:bala}. This is the classical INVERT identity with a weight
$x$ for each part \cite{bernstein}.
Euler's factorial continued fraction is
\[
 F(z)=\cfrac1{1-\cfrac{z}{1-\cfrac{z}{1-\cfrac{2z}{1-\cfrac{2z}{\ddots}}}}}
     =\frac1{1-z/(1-H(z))}.
\]
Successive denominator contraction gives the Jacobi fraction with linear
coefficients $1,3,5,\ldots$ and quadratic numerators
$1^2z^2,2^2z^2,3^2z^2,\ldots$, which is Flajolet's factorial identity
\cite[Theorem 3B, p.~145]{flajolet}. Specifically, the Stieltjes
numerators $s_{2j-1}=s_{2j}=j$ give linear coefficients $s_1=1$ and
$s_{2j}+s_{2j+1}=2j+1$, and quadratic coefficients
$s_{2j-1}s_{2j}=j^2$.
Substitution into \eqref{eq:bala} gives
\[
 G=\frac{1-H-z}{1-H-(1+x)z}
  =1-x+x\frac{1-H-xz}{1-H-(1+x)z},
\]
which is \eqref{eq:balacf}. Every denominator has constant coefficient
one. Continued fractions mean coefficientwise limits of their finite
truncations, so no analytic convergence assertion is involved.
\end{proof}

The input convention matters. Theorem~\ref{thm:main} also applies to
the Catalan input $a_j=C_j$ in A128899, where
$C_j=\binom{2j}{j}/(j+1)$. Replacing $a_j$ by $b_{j-1}$ instead gives
A205574, not the Bell triangle A357583. Similarly,
$a_j=(j-1)!$ gives A084938, not A090238. For the shifted Catalan input
$a_j=C_{j-1}$ it gives the positive rows of A059365, whose displayed
zeroth entry is $0$ rather than our $P_0=1$ \cite{variants}.
These are coefficient substitutions, not claims about the other
DELTA-transform or arithmetic identities in those entries.

\begin{corollary}\label{cor:variants}
For $D(z)=z(1+A(z))$, put $T_D(n,r)=[z^n]D(z)^r$.
Over any unital semiring,
\begin{equation}\label{eq:shift}
 T_D(n+m,m)=\sum_{r=0}^{n}\binom mr T(n,r)\quad(n,m\geq0).
\end{equation}
In particular, this proves the A090238--A084938 binomial identity.
For integer factorial input $a_j=j!$ and $n,r\geq1$,
\begin{equation}\label{eq:factrec}
 rT(n,r)=rT(n-1,r-1)+(n+r-1)T(n-1,r).
\end{equation}
For the Catalan input $a_j=C_j$ and $n,r\geq1$,
\begin{equation}\label{eq:catrec}
 T(n,r)=T(n-1,r-1)+2T(n-1,r)+T(n-1,r+1).
\end{equation}
The last two formulas recover the recurrences recorded in A090238 and
A128899, respectively \cite{factorial,variants}.
\end{corollary}
\begin{proof}
Extract $[z^n]$ in $(1+A)^m=\sum_{r=0}^{m}\binom mr A^r$.
Terms with $r>n$ vanish, giving \eqref{eq:shift}; the binomial theorem
applies because $1$ commutes with $A$.
For factorial input, coefficient comparison gives
$z^2A'=(1-z)A-z$. Thus
$z^2(A^r)'=r(1-z)A^r-rzA^{r-1}$; taking $[z^n]$ gives
\eqref{eq:factrec}. For Catalan input the defining square-root series
in A128899 gives $A=z(1+A)^2$. Multiply by $A^{r-1}$ and extract
$[z^n]$ to obtain \eqref{eq:catrec}. All boundary coefficients come
from $A^0=1$ and $[z^0]A=0$.
\end{proof}

The zero constant term of $A$ is essential for the stated recursion.
For the variant $\widetilde A(z)=1+z^2$ over $\mathbb Z$, define
$\widetilde P_n(x)=\sum_{r=0}^n[z^n]\widetilde A(z)^r x^r$.
The positive input has $a_1=0,a_2=1$, so $xR(2,1)=x$ whereas
$\widetilde P_2=x+2x^2$. This variant is not asserted in the posing
sources; it fails for the displayed example.

\small
\begin{longtable}{@{}>{\raggedright\arraybackslash}p{0.47\textwidth}>{\raggedright\arraybackslash}p{0.46\textwidth}@{}}
\toprule
Source statement & Coverage\\
\midrule
\endhead
Kurkov's prime-row conjecture; arbitrary-sequence assertion
& Proved by Theorem~\ref{thm:main}; includes noncommutative semirings.\\[3pt]
Kurkov's Bell conjecture and arbitrary-sequence proposal (A357583)
& Proved, Corollary~\ref{cor:bell}; same universal theorem.\\[3pt]
Row and weighted sums; alternating sums
& Reproved by Corollary~\ref{cor:sums}; prior identities in A030017,
A030281 and Heinz's 2021 formula in A030018.\\[3pt]
Bell row sums and central terms (A129247, A357584)
& Reproved; prior definitions by Wieder (2008) and Luschny (2022);
Corollary~\ref{cor:bell}.\\[3pt]
Bala's factorial generating function and continued-fraction proposal
(A090238, FORMULA, 27 May 2017)
& Reproved; prior INVERT identity, Bernstein--Sloane (1995), and
factorial fraction, Euler, as proved by Flajolet (1980);
Corollary~\ref{cor:factorial}.\\[3pt]
Diagonal A000079
& $T(n,n)=2^n$ by Corollary~\ref{cor:sums}.\\[3pt]
Columns A000007, A000040, A014342--A014344
& $r=0$ is $\delta_{n,0}$; $r=1$ is the prime sequence for $n>0$;
$r=2,3,4$ are its convolution powers by \eqref{eq:rows}.\\[3pt]
Subdiagonal extraction A340990
& Its defining coefficient $[z^{2n}]A^n=T(2n,n)$ is identified by
\eqref{eq:rows}.\\[3pt]
Convolution triangle A357368, cited for its program
& Theorem~\ref{thm:main} applies to $a_j=j$ for prime $j$ and $a_j=1$
otherwise, the input used in that entry's programs.\\
Bell, factorial and Catalan variants in A357368's comments
& Proved substitutions and binomial shift, Corollary~\ref{cor:variants};
A059365 differs at row zero.\\[3pt]
Factorial and Catalan recurrences (A090238, A128899)
& Reproved source identities, \eqref{eq:factrec} and \eqref{eq:catrec}.\\[3pt]
Row-polynomial programs in A357368; finite-support inputs
& Proved recurrence \eqref{eq:twosided}; the latter is not posed,
but follows for all semirings by truncating its sums.\\[3pt]
Series reversions A007311, A357588; other DELTA-transform identities;
unrelated arithmetic conjectures in prime, power and Bell entries
& Not covered by the shifted-array identity.\\[3pt]
Nonzero constant term in the source series
& Not posed; obstruction proved: $\widetilde A=1+z^2$ fails at $n=2$.\\
\bottomrule
\end{longtable}
\normalsize

\section{Mathematical verification scope}
The commutative-semiring case of \eqref{eq:identity} is formalized by
\texttt{result} in the
\href{https://github.com/the-omega-institute/trureturing/blob/03150b25ffad5ced1bebdca2efc8c05b484c1eca/D5/S1/Recurrence/Algebraic/KurkovConvolutionRowRecursion.lean}{pinned Lean source}
(namespace \texttt{D5.S1.Recurrence.Algebraic.} followed by
\texttt{KurkovConvolutionRowRecursion}). Its private lemma
\texttt{kurkov\_coeff} states
\[
 [x^r]R(n+1,k+1)=[z^n]A(z)^rB_{k+1}(z)
 \quad(n,k,r\in\N)
\]
over every commutative semiring. In that domain the source's product
$x a_kR(n-1,1)$ equals the product in \eqref{eq:recursion}.
The definition \texttt{RecursiveBoundaryRowSeries.tail} in the
\href{https://github.com/the-omega-institute/trureturing/blob/03150b25ffad5ced1bebdca2efc8c05b484c1eca/D5/S1/Recurrence/Algebraic/RecursiveBoundaryRowSeries.lean}{pinned tail-definition file}
has coefficient $a_{k+j}$ in degree $j$ and supplies precisely $B_k$.
The noncommutative extension and the specializations above have the
written proofs in this note. The accompanying \texttt{verify.py}
compares exact convolution coefficients with both the array recursion
and the tail formula in finite integer, Boolean, modular, matrix and
free-symbol examples, the Bell, factorial and shifted triangles, and
the factorial continued fraction; \texttt{finite-checks.json} states
their ranges. These are finite coefficient comparisons, distinct from
the written all-index proofs.

\begin{thebibliography}{9}
\bibitem{a340991}
A. P. Heinz, \emph{Triangle whose columns are self-convolutions of the
primes}, OEIS A340991, revision 24, 6 April 2025; FORMULA conjecture by
M. Kurkov, 5 April 2025.
\url{https://oeis.org/history/view?seq=A340991&v=24}.

\bibitem{osullivan}
C. O'Sullivan, \emph{De Moivre and Bell polynomials},
arXiv:2203.02868v1, 2022, Definition 1.1 and \S2.
\url{https://arxiv.org/abs/2203.02868v1}.

\bibitem{bell}
P. Luschny, \emph{Convolution triangle of the Bell numbers}, OEIS
\href{https://oeis.org/history/view?seq=A357583&v=10}{A357583, revision 10},
6 April 2025; FORMULA conjecture by M. Kurkov, 5 April 2025.
Related entries: T. Wieder, \href{https://oeis.org/A129247}{A129247}
(2008); P. Luschny, \href{https://oeis.org/A357584}{A357584} (2022).

\bibitem{factorial}
P. Del\'eham, \emph{Lists of unlabeled permutations}, OEIS
\href{https://oeis.org/history/view?seq=A090238&v=46}{A090238, revision 46},
17 October 2022; title from a comment of D. Callan; FORMULA generating
function and continued fraction by P. Bala, 27 May 2017.

\bibitem{bernstein}
M. Bernstein and N. J. A. Sloane, \emph{Some canonical sequences of
integers}, Linear Algebra Appl. 226--228 (1995), 57--72;
erratum 320 (2000), 210. \S1, INVERT.
\url{https://arxiv.org/abs/math/0205301}.

\bibitem{flajolet}
P. Flajolet, \emph{Combinatorial aspects of continued fractions},
Discrete Math. 32 (1980), 125--161, Theorem 3B, p.~145.
\url{https://algo.inria.fr/flajolet/Publications/Flajolet80b.pdf}.

\bibitem{variants}
The OEIS Foundation, convolution-triangle variants
\href{https://oeis.org/A205574}{A205574},
\href{https://oeis.org/A084938}{A084938},
\href{https://oeis.org/A128899}{A128899},
\href{https://oeis.org/A059365}{A059365}; reversion entries
\href{https://oeis.org/A007311}{A007311},
\href{https://oeis.org/A357588}{A357588}; accessed 13 November 2026.

\bibitem{sequences}
The OEIS Foundation, entries
\href{https://oeis.org/A030017}{A030017},
\href{https://oeis.org/A030018}{A030018},
\href{https://oeis.org/A030281}{A030281},
\href{https://oeis.org/A340990}{A340990} and
\href{https://oeis.org/A357368}{A357368};
sequence definitions and formulas, accessed 13 November 2026.

\bibitem{transforms}
N. J. A. Sloane, \emph{Transformations of integer sequences}, INVERT and COMPOSE entries,
The On-Line Encyclopedia of Integer Sequences.
\url{https://oeis.org/transforms.txt}.
\end{thebibliography}

\bigskip\noindent\small
\copyright\ 2026 Haobo Ma and Wenlin Zhang.
Licensed under \href{https://creativecommons.org/licenses/by/4.0/}{Creative
Commons Attribution 4.0 International (CC BY 4.0)}.
\end{document}
